% Chapter V -- Results and Discussion (rev L-104)
% Interpretive, not enumerative: numbers live in tables regenerated from
% run artifacts; the chapter argues what they MEAN. Every numeric slot is
% an explicit [[FILL]] token bound to a named artifact -- transcription
% replaces invention.
% REV L-104: fixed stray "#" comment character (compile error); softened
% "converts to confirmatory" overclaim; added [[VERIFY-AGAINST]] gates on
% directional claims; inserted [CITATION NEEDED] markers on uncited
% literature assertions.

\chapter{Presentation, Analysis, and Interpretation of Results}

\section{Reading Guide}
Results are organized around the four research questions. For each, we
report the estimate, its paired uncertainty, the audit that stress-tested
it, and the interpretation. The pilot study is presented first as
instrumentation evidence, not as confirmatory result: its role was to
verify pipeline mechanics, surface coding errors, and fix the analysis
script before unblinding the full study.
[[VERIFY-AGAINST: every directional claim in this chapter must match
full\_study\_results.md / paired-delta artifacts exactly once FILL tokens
are transcribed; delete this token at that time.]]

\section{Pilot Study: Instrumentation Findings}
The pilot executed the end-to-end path---generation, guarding, training,
scoring, auditing---at reduced scale. Outcomes ([[FILL: implementation/
results/pilot\_results.md]]): pipeline ran to completion without manual
repair; early metric implementations agreed with reference computations;
preliminary AUROC ordering matched expectations, supporting the decision
to proceed to full scale without protocol modification; two
instrumentation defects (binning edge cases in ECE, a stratum-label
collision) were found and fixed, with fixes locked before full-study
execution. \emph{Interpretation}: freezing the corrected analysis code
before full-study execution removes the most common exploratory biases;
the full study is accordingly treated as \emph{pre-specified}. This is a
weaker claim than ``confirmatory,'' which would require independent
replication; the distinction is carried into the limitations discussion.

\section{RQ1 --- Surrogate-Trained Detection on Held-Out Tiers}
Table~\ref{tab:rq1-headline} reports headline AUROC/AUPRC with 95\%
bootstrap intervals by data tier (external RAID strata; own-corpus test
split; dual-reported variants excluding recovered shards).
[[FILL: full\_study\_results.md headline table; reproduce rows verbatim.]]

\emph{Analysis}. Three patterns organize the reading, contingent on the
transcribed numbers above. First, usable-but-degraded transfer: surrogate-
trained detection is expected to retain discriminative signal on real text
while falling below in-distribution ceilings, quantifying the
surrogate--real gap rather than assuming it away. Second, tier asymmetry:
any difference between Tier F1 and Tier E performance localizes whether
the gap is generator-driven or collection-driven --- collection-driven
gaps implicate acquisition protocol, generator-driven gaps implicate
training distribution. Third, shard-insensitivity: the with/without-F2
contrast bounds the influence of the salvage pathway on headline claims
(cross-ref Section~6.3).

\section{RQ2 --- Paired Comparison Against Baselines}
Figure/Table pair [[FILL: paired-deltas artifact]] gives per-stratum
paired deltas of the principal system against zero-shot statistical and
embedding baselines, with paired-bootstrap intervals and Holm-adjusted
significance stars over the prespecified comparison family.
[[FILL: delta table verbatim.]]
\emph{Analysis}. Deltas are read jointly with intervals, never as point
verdicts. Where the surrogate-trained system wins uniformly, training on
synthetic proxies is vindicated despite the gap of RQ1; where wins invert
on specific domains, the inversion pattern feeds directly into the
headroom decomposition (Section~5.6). Baseline parity on any stratum is
itself informative: it caps the marginal value of surrogate training there
and redirects effort toward data acquisition rather than modeling.

\section{RQ3 --- Calibration and Reliability Across Strata}
Reliability diagrams per stratum and binned ECE [[FILL: reliability curve
artifacts]] answer whether the detector's probabilities deserve trust, a
question orthogonal to ranking accuracy.
\emph{Analysis}. The failure mode discussed in the literature---over-
confident misranking under distribution shift [CITATION NEEDED: wire the
specific calibration-under-shift source(s) during Chapter~2 migration]---
is tested against the observed direction [[FILL: over-/under-confidence,
worst strata]]. Practically, a detector that ranks well but calibrates
poorly supports triage (flagging) but not threshold automation; this
distinction disciplines the recommendation language in Chapter~6 away
from deployment overclaim.

\section{RQ4 --- Where the Error Lives: Headroom, Length, Authorship}
The headroom-by-length audit decomposes residual error into length bands
([[FILL: headroom\_by\_length artifact]]), the mixed-authorship breakdown
isolates contamination-sensitive slices, and paraphrase/domain breakdowns
plus repair-attempt logs complete the localization.
\emph{Analysis}. Aggregate AUROC conceals slice pathology wherever slices
disagree with the pooled number [[FILL: e.g., short-text band retains
largest reference-gap; mixed-authorship stratum behaves as neither
class]]. Repair attempts against the largest localized gaps are reported
even where unsuccessful, converting negative results into boundary
knowledge: they delimit what current surrogate-training recipes cannot
yet buy.

\section{Synthesis Against Prior Literature}
Chapter~2 positions this study against three literatures: statistical
detection (likelihood/perplexity tradition), representation-based
detection, and robustness/evasion analyses. The present results speak to
each, conditionally on the transcribed numbers: (a) if the paired deltas
of Section~5.4 hold as tabulated, the surrogate-training strategy achieves
transfer that zero-shot statistics do not, strengthening the benchmark-
construction line of the RRL relative to the pure-statistics line;
(b) calibration degradation under shift, if observed as templated in
Section~5.5, corroborates critiques that headline AUROC benchmarks
overstate deployment readiness [CITATION NEEDED: attribute the critique,
e.g., to the RAID robustness evaluation and successor work]; (c) evasion
sensitivity quantifies, on controlled strata, the fragility arguments made
in the paraphrase-evasion literature [CITATION NEEDED: bind to
krishna2023paraphrase at migration]. The novelty chain from the idea-phase
work---intrinsic-dimension probes, linguistic splitting as a profiling
lens, evader suites---lands here as the audit apparatus that makes these
claims slice-resolved rather than aggregate, which is the study's
principal methodological contribution beyond its empirical ones.

\section{Summary of Chapter Findings}
(1) Surrogate-only training yields cross-tier detection whose usability
and transfer cost are quantified in Sections~5.3 [[VERIFY-AGAINST: confirm
gap magnitude from E-02]]. (2) Paired deltas establish superiority over
zero-shot baselines on [[FILL: count]] of prespecified strata with
interval-supported margins. (3) Calibration degrades under shift in
direction [[FILL]], constraining deployment claims to triage. (4) Residual
error concentrates in [[FILL: length band / authorship / attack stratum]],
partially repaired by [[FILL: attempt]], with remaining headroom owned
explicitly in Chapter~6.